\documentclass[12pt]{article}
\usepackage[margin=1.0in]{geometry}
\usepackage{setspace}
\setstretch{2}

\usepackage{hyperref}
\hypersetup{%
    colorlinks=true,
    linkcolor=cyan,
    filecolor=magenta,      
    urlcolor=blue,
    citecolor=blue,
}
\usepackage{siunitx}
\DeclareSIUnit{\molar}{M} 

% Bibliography
\usepackage{natbib}
\bibliographystyle{msb3} % the optional note field in the bib file is useful for MSB data refs

\begin{document}

\begin{titlepage}
    \begin{center}
        \vspace*{1cm}

        Atomic elementary flux modes explain the steady state flow of
        metabolites in large-scale flux networks

        \vspace{1.0cm}
        aEFMs explain metabolite flows
            
        \vspace{1.0cm}

        Justin G. Chitpin$^{1,2}$ and Theodore J. Perkins$^{1,2*}$

        \vspace{1.0cm}

        $^1$Ottawa Hospital Research Institute, 501 Smyth Road, Ottawa, K1H 8L6, Ontario, Canada\\
        $^2$Ottawa Institute of Systems Biology, Department of Biochemistry, Microbiology and Immunology, University of Ottawa, 451 Smyth Road, Ottawa, K1H 8M5, Ontario, Canada
           
        \vspace{0.5cm}

        $^*$ Corresponding author.

        \vspace{1.0cm}
     
        Flux network/Elementary flux mode/flow decomposition/Markov chain/atom mapping
            
   \end{center}
\end{titlepage}

\clearpage

\section*{Abstract} % 175 words max

Steady state fluxes are a measure of cellular activity under metabolic
homoeostasis, but understanding how individual substrates are metabolized
remains a challenge in large-scale networks. Pathway-based approaches such
as elementary flux mode (EFM) analysis are limited to small networks due
to the combinatorial explosion of pathways and ambiguity decomposing
fluxes onto EFMs. Here, we present an alternative approach to explain
metabolic fluxes in terms of the steady state flow of their atomic 
constituents. We refer to these pathways as atomic EFMs (AEFMs) and show
that computations involving AEFMs are orders of magnitude faster than
standard EFMs which balance molecular stoichiometries. Using our approach,
we enumerate carbon and nitrogen AEFMs in five genome-scale metabolic
models and compute the AEFM decomposition of fluxes estimated in a HepG2
liver cancer cell line. Our results systematically characterize carbon and
nitrogen remodelling and, on the HepG2 network, predicts major
anapleurotic pathways of glutamine carbon metabolism, including the
recently discovered non-canonical TCA cycle. We anticipate our method will
be useful for quantifying cellular metabolism as high-throughput fluxomic
analyses become increasingly common.

\clearpage

\section*{Introduction} % intros tend to be very succinct

% (a) flux overview
Recent technological and computational advancements are leading to a rise in
metabolic flux
data~\citep{LeeWonDong2019Spas,WangLin2022Srit,WangRuohong2022Gstm}. These
fluxes quantify the rate of metabolite interconversions within a cell or their
movements across subcellular compartments. Metabolic fluxes are studied to
understand the dynamics of substrate
switching~\citep{ZampieriMattia2012Rmuc,KochanowskiKarl2021Gcom}, redox
balancing~\citep{McKinlayJamesB2010Cdfa,FranchinaDavideG2022Grbc}, overflow
metabolism~\citep{BasanMarkus2015OmiE,VanderHeidenMatthewG2009UtWE}, cell
differentiation~\citep{EndoYusuke2022Fami,BenjaminJacksonT2024Mrot}, and
cellular
communities~\citep{SchuetzRobert2012MOoM,YuJasonSL2022Mcfr,Gustafsson2024Mcbc}.
Hence, `fluxomics'~\citep{EmwasAbdul-Hamid2022F-NM,SanfordKarl2002Gtfa} is
emerging as a field dedicated to quantifying the totality of fluxes within
biological systems.

% (b) estimating steady state fluxes - happy with the writing
Current experimental and computational methods generally estimate fluxes
under metabolic homeostasis. Under this assumption, the production and
consumption of metabolites are balanced over time, and the fluxes are said
to be under steady state. Techniques such stable isotope tracing analysis
(SITA) experimentally determine steady state fluxes by culturing cells in
labelled substrates and quantifying their incorporation within endogenous
metabolites~\citep{BuescherJoergM2015Arfi,AntoniewiczMaciekR2018Agt1}.
Using this method, fluxes have been measured in major metabolic subsystems
such as glycolysis, pentose phosphate pathway, and citric acid
cycle~\citep{OrthJeffreyD2010RaUo,AhnEunyong2017Tfro,WuChao2019FARC}. Many
computational methods have also been proposed to infer steady state fluxes
from transcript, protein, and metabolite abundance
data~\citep{UematsuSaori2022Mlmf,GonzalezArrueNicolas2023Peom,HuangYuefan2023Ccmf,KasteJoshuaAM2023Afpu,LeeJustinY2023Tiac}.
Collectively, these approaches are being used to estimate and predict
fluxes in genome-scale metabolic models (GEMs) involving hundreds to
thousands of
reactions~\citep{NilssonAvlant2020Qaoa,CherkaouiSarah2022Afao,GelbachPatrickE2024Fsig}.

% (c) introducing problem
Although much effort has been made to estimate steady state fluxes, less
emphasis has been placed on developing methods to analyze these
large-scale flux datasets. In contrast to abundance data (e.g.
transcripts, proteins, metabolites), steady state fluxes are constrained
by the reaction stoichiometries between metabolite inflows and outflows.
In a network with fixed source and sink fluxes an increase in flux along
one metabolic pathway must decrease fluxes along one or more other
pathways to maintain mass conservation. For simple networks, one may
predict how the system would evolve by identifying the up and downstream
reactions flanking an enzymatic perturbation. For example,
\citeauthor{SchwartzJean-Marc2006Qema} were able to enumerate all
glycolytic pathway fluxes in a 12 reaction model of yeast
glycolysis~\citep{SchwartzJean-Marc2006Qema}. Yet in more complex genome-scale
metabolic networks containing thousands of fluxes
(e.g.~\cite{BrunkElizabeth2018Reat,RobinsonJonathanL2020Aaoh,WangHao2021Gmnr}),
there exists no hierarchical relationship between reactions. A change in one
reaction rate, let alone multiple perturbations, may alter network fluxes
beyond neighbouring reactions~\citep{NobileMarcoS2021Agsa}.

% (d) EFMs
This problem of understanding the flow of metabolites within steady state
flux networks has historically been addressed through the concept of
elementary flux modes (EFMs). EFMs are a pathway-based approach to analyze
metabolic networks and are defined as minimal sets of biochemical
reactions that maintain steady state flux~\citep{SchusterStefan1994Oefm}.
This minimal property specifies that an EFM cannot be decomposed into two
or more smaller pathways carrying steady state flux. An attractive
property of EFMs is that any set of steady state fluxes in a network can
be explained as a positive, linear combination of
EFMs~\citep{SchusterStefan1999Doef}. While EFMs are a mechanistic
explanation of how metabolites travel through networks, the number of EFMs
scales combinatorially as a function of network
size~\citep{AcunaVicente2009Maci}. It remains computationally infeasible
to enumerate EFMs in GEMs after over a decade of algorithmic advancements
(e.g.
\citealt{GagneurJulien2004Coem,TerzerMarco2008Lcoe,BuchnerBiancaA2021EaPp}).
The number of EFMs also poses a challenge for explaining network fluxes in
terms of their EFMs. When there are more EFMs than linearly independent
fluxes, there exist multiple ways to assign EFM weights that reconstruct
the observed network fluxes. Consequently, structural analyses of EFMs and
downstream analyses of their weights have been limited to small-scale
metabolic
subnetworks~\citep{FakihIbrahim2023Dgmm,MpabanziLiliane2022Frca,ToyaYoshihiro2022Mpef,MatsudaFumio2011Esoy,SchwartzJean-Marc2006Qema}.

In our previous work, we took the first steps to solve the EFM flux
decomposition problem for the special case of unimolecular reaction
networks under a Markov constraint. We modelled a hypothetical particle
moving between metabolites under the assumption that the transition
probability of each reaction was proportional to the observed steady state
flux. By expanding the notion of Markov chain state to include sequences
of metabolites, we proposed an algorithm to efficiently compute simple
cycle probabilities corresponding to a given EFM. We called this construct
the \emph{cycle-history Markov chain} (CHMC) and proved that these EFM
probabilities could be scaled to weights that fully reconstructed the flux
along each reaction~\citep{ChitpinJustinG2023AMct}.

% (e) directives
In this work, we generalize our CHMC method to operate on any type of
metabolic model including those involving multispecies reactions. We do so
by proposing a new type of \emph{atomic elementary flux mode} which
traces the steady state flow of an individual atom within a metabolic
network. We refer to these as AEFMs to distinguish them from the
(molecular) EFMs proposed by \citeauthor{SchusterStefan1994Oefm}. While
a linear combination of molecular EFMs will reconstruct the overall reaction
fluxes, the totality of AEFM weights for a given source metabolite will always
reconstruct its mass flow entering the network. We first introduce our
framework for enumerating and uniquely identifying AEFM weights. Using
a state-of-the-art atom mapping algorithm, we show how to generate atomic
transition graphs from which we construct atomic cycle-history Markov chains
(ACHMCs). From these ACHMCs, we enumerate AEFMs in GEMs considered infeasible
by state-of-the-art molecular EFM enumeration programs. We further analyse the
structure of the AEFMs to characterize nutrient source remodelling. Lastly, we
analyze a single set of AEFM weights estimated from a HepG2 liver cancer cell
line to quantify glutamine metabolism. Our results predict glutamine in HepG2
cells is metabolized through a non-canonical TCA pathway recently discovered in
non-small cell lung cancer (NSCLC) cells and C2C12 mouse
myoblasts~\citep{ArnoldPaigeK2022Anta}.

\section*{Results}

\subsection*{Tracing atomic flows in genome-scale models} % pipeline overview

We first describe what molecular EFMs look like in metabolic networks and
provide intuition why enumerating and decomposing fluxes onto these
pathways remains a fundamental problem in literature. Figure 1a shows an
example molecular EFM involving substrates glucose and ATP which are
consumed to produce metabolites G3P, DHAP, and ADP. Multiple source and
sink fluxes are required to balance the internal metabolite
stoichiometries, leading to several branching and converging paths within
the molecular EFM. As the metabolic networks expands to include dozens or
more reactions, the corresponding molecular EFMs can grow increasingly
complex with numerous source and sink fluxes required to
stoichiometrically balance participating metabolites. In particular, there
may be multiple sets of reactions involved to balance pairs of energetic
molecules (e.g. ATP/ADP), redox equivalents (e.g. NAD$^+$/NADH), metabolic
byproducts (e.g. water, hydrogen ions), or coenzymes (e.g. coenzyme A).
This complexity can lead to an intractable number of pathways that may
only differ by reaction subsets required to mass-balance these
metabolites. Choosing metabolic byproducts to remove may, in part, solve
this problem. However, the choice of metabolites to prune from the network
is often subjective and not guaranteed to sufficiently reduce the total
number of molecular EFMs.

Given many metabolic flux networks are generated to study nutrient
catabolism~\citep{ZamboniNicola2015DtMS}, we argue for studying source
metabolite-derived AEFMs to reduce the combinatorial number of molecular
EFMs. An example carbon AEFM is highlighted in Figure 1b which shows how
the C6 glucose carbon is remodelled during the first steps of glycolysis.
As atoms are never split apart or fused together in metabolic reactions,
these AEFMs always correspond to simple paths spanning source and sink
metabolites or simple cycles between internal metabolites. These pathways
are straightforward to interpret because each atom within an AEFM can only
transit from a single substrate to product, regardless of the network
size, connectivity, or number of participating substrates/products in each
reaction. This property further enables us to enumerate and compute the
AEFM weights using our previously described CHMC method which can uniquely
decompose fluxes onto EFMs for unimolecular reaction
networks~\citep{ChitpinJustinG2023AMct}.

We outline our computational pipeline in Figure 1c to enumerate AEFMs and
uniquely assign their weights when steady state fluxes of the network are
known. In step 1, we require as inputs the metabolic network
defined by a set of metabolites and reactions encoded as a stoichiometry
matrix, the metabolite structures denoted as SMILES strings, and the
steady state fluxes. From the stoichiometry matrix and SMILES strings, we
construct reaction SMILES strings which are imputed into the atom mapping
program RXNMapper~\citep{Schwaller2021extraction} to map the position of
each atom within each reaction equation. This atom mapping data is used to
construct an atomic transition graph in step 2, which traces the atomic
position of an individual atom from a given source metabolite entering the
network. We set the transition probabilities proportional to the atomic
flux through each reaction, accounting for those with multiple substrate
stoichiometries.
%For example, 10\SI{}{\milli\molar} of input glucose flux would correspond
%with 10\SI{}{\milli\molar} of atomic flux through each glucose carbon.
From each atomic transition graph, we construct the corresponding ACHMC to
enumerate all AEFMs involving that source metabolite atom in step 3.
Finally, the AEFM weights are computed in step 4 by steady state analysis
of the ACHMC. 

\subsection*{Enumerating AEFMs is computationally tractable in large-scale networks}

Given the major computational challenge of enumerating molecular EFMs in
networks with $>100$ metabolites and reactions, we first sought to test
whether it was feasible to enumerate atomic EFMs in GEMs. We selected five
genome-scale metabolic models across distinct organisms in ascending order
of network size. These networks include a metabolic model of \textit{E.
coli}, human red blood cells, \textit{H. pylori}, \textit{S. aureus}, and
HepG2 liver cancer cells. These models were pre-processed to ensure
unambiguous atom mappings across all reactions. Namely, all metabolites
with no known structures (pseudometabolites), and pseudoreactions with
non-integer stoichiometries were removed from the networks. The resulting
metabolic networks ranged between 71-547 metabolites and 144-974 reactions
(Figure EV1). We then attempted to enumerate the molecular EFMs using
FluxModeCalculator which is the fastest molecular EFM enumeration
program~\citep{vanKlinkenJanBert2016Faet}. Following our pipeline, we
subsequently compiled the metabolite SMILES strings and constructed ACHMCs
rooted on each carbon and nitrogen atom across all source metabolites. We
focused on carbon and nitrogen atoms, in particular, for their metabolic
significance over chemical elements such as oxygen or phosphorus. The
resulting summary statistics describing the GEMs and ACHMCs are presented
in Table 1. Across all carbon ACHMCs, we observed a broad increase in the
total number of AEFMs as a function of metabolic network size. The one
exception we observed was the HepG2 network which contained fewer
metabolites and reactions than iSB619, yet exhibited two orders of
magnitude more carbon AEFMs, suggesting a greater degree of network
connectivity. 
%Collapsing across carbon AEFMs and nitrogen AEFMs, we observed
%a consistent ratio between the number of ACHMC transitions and ACHMC
%states ($\mu_C = 1.43, \sigma^2_C = 0.14$; $\mu_N = 1.53, \sigma^2_N
%= 0.16$).

Figure 2a shows the total number of atomic and molecular EFMs enumerated
by our program MarkovWeightedEFMs.jl and FluxModeCalculator on the five
GEMs. Using our ACHMC pipeline, we enumerated between $10^3 - 10^6$ carbon
and nitrogen EFMs within each GEM. In contrast, FluxModeCalculator failed
to enumerate molecular EFMs in all but the two smallest E. coli core and
iAB RBC 283 datasets, returning $10^3 - 10^4$ times more
molecular-versus-atomic EFMs. By focusing on an atom of interest, our
ACHMC approach greatly reduced the number of EFMs to biologically relevant
pathways involving nutrient source remodelling. This reduction in pathways
resulted in AEFM-versus-EFM enumeration completing $10^3-10^4$ times
faster (Figure 2b). While running FluxModeCalculator in parallel did
improve run times by 8.2-8.6 times (Figure EV2), it remained over
$10^1-10^2$ slower than enumerating the corresponding carbon and nitrogen
AEFMs of each GEM. We further found that AEFM enumeration actually
involved fewer metabolite-atom combinations than the total number of
metabolites in a given network. Constrained by the atom mapping
predictions, we observed that the average carbon and nitrogen ACHMC across
all five GEMs only transited 0.41\% and 1.7\% of the total number of
metabolite-carbon and metabolite-nitrogen states, respectively (Figure
EV3). While serial run times in Figure 2c scaled quadratically (Figure 2c;
$\text{slope}=1.90, \text{intercept}=-6.54$), our program and AEFM
approach benefits greatly from parallelism. Each ACHMC can be computed
independently of one another and our program is also parallized on the
level of individual ACHMCs to improve AEFM enumeration. For the top 28
largest carbon ACHMCs in the HepG2 network, Figure 2d shows a median
speedup of 7.1 times when using 32 threads. 


\iffalse

These fast AEFM enumeration times may seem counterintuitive given the
number of metabolite-atom positions is much larger than the total number
of metabolites in a given metabolic network. To resolve these findings, we
enumerated the total number of metabolite-carbon and metabolite-nitrogen
positions across all five GEMs. The HepG2 and iAB RBC 283 networks
exhibited the greatest number of combinations with 8194 metabolite-carbon
and 2533 metabolite-nitrogen positions. This corresponded to over 18 and
4 times the number of metabolites within each respective network. However,
  each atomic transition graph may only involve a subset of the entire
  list of metabolite-carbon or metabolite-nitrogen positions. Constrained
  by the atom mapping predictions, we found that the average carbon and
  nitrogen ACHMC across all five GEMs only transited 0.41\% and
1.7\% of the total number of metabolite-carbon and metabolite-nitrogen
states, respectively (Figure EV3). Hence, the metabolite-atom state space
of any given ACHMC is several orders of magnitude smaller than the total
number of network metabolites.

We therefore hypothesized that the speedup between MarkovWeightedEFMs.jl
and FluxModeCalculator was due to the proportionally fewer number of
AEFMs-versus-molecular EFMs across the metabolic networks. For molecular
EFM enumeration algorithms, it is well established that run times cannot
be predicted from the number of metabolites or reactions within the
model~\citep{UllahEhsan2020Tsef}. We found, however, that AEFM run times
scale approximately quadratically as a function of the top 28 largest
ACHMCs (Figure 2c; $\text{slope}=1.90, \text{intercept}=-6.54$). While
future work is required to improve the algorithm scaling, we emphasize
that the computation of each ACHMC is embarrassingly parallel at the
source metabolite-atom level. Our software is also parallelizable on the
level of individual ACHMCs to further speed up AEFM enumeration. For those
top 28 largest carbon ACHMCs in the HepG2 network, we re-enumerated their
AEFMs with varying numbers of threads. Figure 2d reports the algorithm
scaling, where we observed a median speedup of 7.1 times when using 32
threads. As expected, the greatest speedups were associated with the
largest ACHMCs. For example, the largest ACHMC corresponded to a glutamine
carbon with 1,120,290 ACHMC states/168,793 AEFMs, which exhibited a 9.3
times speedup when enumerated with 32 threads.

\fi

\subsection*{Most AEFMs are source-to-sink pathways spanning dozens of reactions}

We next turned to characterizing the structural properties of all carbon
and nitrogen AEFM across the five GEMs. Molecular EFMs can be divided into
pathway that span source and sink metabolites or internal loops within the
metabolic network such as pairs of reversible reactions. AEFMs can be
further classified into those that transit sets of distinct metabolites or
those that revisit the same metabolite, albeit in different atomic
positions. An example of source-to-sink AEFM with this characteristic is
a glucose-derived carbon AEFM remodelling through two turns of the TCA
cycle; on the second turn, the position of the carbon atom shifts within
the TCA intermediates, resulting in different isotopomers and therefore
a source-to-sink pathway with a ``metabolite revisitation'' under our
definition.

To obtain a broad understanding of the types of AEFMs within the five
GEMs, we categorized them into source-to-sink pathways and internal cycles
with or without metabolite revisitations in different atomic positions.
Figure 3a-b shows the counts for each AEFM classes within each GEM. We
observed that the majority of carbon AEFMs were source-to-sink pathways
with or without metabolite revisitations. Those without metabolite
revisitations tended to scale as a function of metabolic network size,
while the number of looped carbon AEFMs with or without metabolite
revisitations remained similar across the datasets. Investigation of
source-to-sink pathways confirmed that many metabolite revisitations
involved the TCA cycle intermediates. This explained why the iAB RBC 283
dataset contained the fewest number of these pathways since they lack
cellular organelles and cannot support cyclic reactions associated with
metabolite transport between compartments or mitochondrial TCA activity.
Biological examples of internally looped AEFMs with metabolite
revisitations varied across the GEMs, but generally consisted of lengthy
pathways also involving TCA metabolites. The greatest number of these
pathways were found in the E. coli core network, with many metabolite
revisitations involving carbon flux through the glyoxylate shunt or carbon
dioxide, from TCA cycle activity, re-entering central carbon metabolism
through the production of oxaloacetate. Similar findings were observed for
the nitrogen AEFMs, although looped pathways with metabolite revisitations
were only observed in the iIT341 and iSB619 networks through reactions
involving the production and consumption of ammonia. 

Figure 3c-d shows the distribution of atomic EFM lengths across the five
GEMs. The carbon AEFMs from all but the HepG2 networks tended to be quite
short with few pathways exceeding 40 reactions. In contrast, the HepG2
carbon AEFM lengths exhibited a Gaussian distribution with mean of 61.7
and variance of 15.5. As there is a greater number of metabolites and
reactions in the iSB619 network, these results suggest a higher degree of
connectivity between metabolites within the HepG2 network. In particular,
to emphasize how different the paths available to each atom in
a metabolite may be, we computed the number of AEFMs for each carbon atom
(resp. nitrogen atom). For each metabolite, we then computed the ratio of
the maximum number of paths for any carbon (resp. nitrogen), and the
minimum number of paths. So for example, a ratio of 10 means that one of
the carbon atoms can transit the network in 10 times as many different
ways as the carbon with the fewest paths through the network. Our results
in Figure 3e-f shows that a surprising number of carbons and nitrogens can
transit different pathways within their respective source metabolites.
Specifically, many hydrophobic and positively charged amino acids such as
phenylalanine, methionine, arginine, and histidine could undergo carbon
and nitrogen ratios over $10^2$. Source metabolites containing fewer
carbons or nitrogens tended to exhibit ratios closer to 1, such as glycine
which is equivalently metabolized across the same AEFMs in each GEM.


% More comments?

% 209 source mets with non-zero carbon statistics
% 443 source mets total with carbons
% 47.2 percent 
% 16 source mets with only one carbon

% 70 source mets with non-zero nitrogen statistics
% 202 source mets total with nitrogen
% 34.6 percent  
% 108 source mets with only one nitrogen

\subsection*{A minority of carbon AEFMs explains the majority of source metabolite remodelling}

% Glutamine carbon 1:   8,098
% Glutamine carbon 2: 168,793
% Glutamine carbon 3:  42,691
% Glutamine carbon 4:  16,581
% Glutamine carbon 5:   4,922
% Glutamine carbon AEFM total: 241,085
% Glutamine carbons can remodel into 202 distinct metabolites aside from
% input glutamine

We finally highlight the utility of our AEFM method by analyzing a single set
of steady state fluxes in the HepG2 dataset. Specifically, we focused on
identifying alternative pathways of glutamine metabolism since many cancer
cells, including HepG2, are dependent on this non-essential amino acid for
survival and proliferation~\citep{YeYanyan2023Gmri,YooHeeChan2020Gric}.
Using our method, we decomposed the set of steady state fluxes onto the
241,085 glutamine carbon AEFMs characterizing the HepG2 metabolic network.
Figure 4a shows the cumulative explained mass flow of each glutamine
carbon as a function of their AEFMs ranked from greatest to smallest
explained mass flow. We capped the cumulative explained mass flow of each
glutamine carbon at 99\% because we observed that hundreds to thousands of
carbon AEFMs collectively explained the remaining 1\% glutamine mass flow.
The resulting curves indicate that the majority of glutamine carbon
remodelling is explained by a surprisingly small number of glutamine
AEFMs, especially since our Markovian flux decomposition assigns some
weight to all network pathways. The top 5 ranked AEFMs across each
glutamine carbon (25 pathways altogether) explained over half (50.2\%) of
the total glutamine carbon mass flow with the subsequent
5 highest ranked AEFMs only explaining an additional 9.8\% glutamine
carbon mass flow. These observations were consistent across carbon AEFMs
derived from other amino acids as well as glucose, suggesting
a \emph{dominant flux hypothesis} where steady state pathway fluxes are
centralized to a subset of AEFMs (Figure EV3). For some source
metabolites, we observed overlapping curves between carbons (e.g.
arginine, glycine, serine), suggesting individual carbons were remodelled
through the same AEFMs. Corroborating results from Figure 3e-f, we
observed many non-overlapping curves for the majority of amino acids since
their carbons could transit the network via different internal
metabolites.

%We denote this carbon mass flow as $C_k = w_k \cdot L_k$, where $L_k$ is
%the length of EFM $k$. 

To obtain a high-level understanding of the dominant pathways of glutamine
carbon remodelling, we constructed a metabolic subnetwork from the top
5 glutamine carbon AEFMs, collapsing metabolite-atom positions onto distinct
metabolites. Of the 435 metabolites in the HepG2 dataset, input glutamine
can remodel into 202 species localized to the cytosol or mitochondria. We
found that the top 5 glutamine-derived carbon AEFMs consist of just 23
metabolites and 31 reactions (Figure 4b). These AEFMs generally consisted
of cyclic pathways involving TCA intermediates. Only two source-to-sink
reactions converting cytosolic glutamine towards the protein pool and
carbon dioxide through alpha-ketoglutarate decarboxylation in the
mitochondria. 

As AEFMs are steady state pathways, they can be visualized as Sankey
diagrams to better understand the carbon flows through the network. Figure
4c-e shows the top 5 AEFMs for each glutamine carbon with metabolites
represented as nodes and each AEFM as a flow proportional to its weight.
Although each glutamine carbon can transit a different number of AEFMs, we
found the top 5 backbone carbons and first two sidechain carbons were
identical in both AEFM pathway and weight.
    
Focusing on either backbone carbons in Figure 4c, we observed that all
AEFMs corresponded to well-established metabolic subsystems. The first,
third, and fifth highest ranked AEFMs corresponded to metabolites and
reactions involved in the malate-aspartate shuttle. The fourth largest
AEFM closely resembled canonical TCA cycle activity, albeit with TCA
intermediates citrate, isocitrate, and alpha-ketoglutarate localized to
the cytosol rather than mitochondria. This AEFM also involved additional
reactions transforming alpha-ketoglutarate into glutamate before
remodelling into succinyl-CoA in the mitochondria. We noticed that the
second largest AEFM cycled carbon mass between TCA intermediates citrate,
oxaloacetate, and malate between the cytosol and mitochondria.
Interestingly, this metabolic pathway corresponds to the recently
discovered non-canonical TCA activity described by
\citeauthor{ArnoldPaigeK2022Anta}, who validated this pathway in non-small
cell lung cancer cells, C2C12 mouse myoblasts, and mouse embryonic stem
cells by stable isotope tracing analysis~\citep{ArnoldPaigeK2022Anta}.
This non-canonical TCA activity was also observed in glutamine carbons
3 and 4 in AEFM 2 (Figure 4e). We also found AEFMs 1, 3, and 5 in Figure
  4d-e also corresponded to the malate-asparate shuttle. Only the terminal
  glutamine carbon in Figure 4e contained source-to-sink AEFMs that
  decarboxylated glutamine-derived alpha-ketoglutarate into carbon
  dioxide.

%Having mapped individual Sankey diagram AEFMs to existing metabolic
%subsystems, we asked whether the results could be combined into a 

Our analysis of the most active AEFM weights revealed pathway activity in
well-established metabolic subsystems. This finding led us to ask whether
we could summarize our AEFM analysis as a simplified metabolic model
describing the majority of glutamine carbon remodelling in the HepG2 flux
network. The resulting schematic is shown in Figure 4f which consists of
the following three major metabolic pathways: canonical TCA cycle,
non-canonical TCA cycle, and the malate-aspartate shuttle. These metabolic
pathways were chosen based on the number of top glutamine-derived carbon
AEFMs that overlapped with reactions within these subsystems. Our results
suggest input glutamine is converted into cytosolic glutamate which is
transported into the mitochondria through aspartate/glutamate antiporters
encoded by \textit{SLC25A12} and \textit{SLC25A13}. Mitochondrial
glutamate is then converted into TCA intermediate alpha-ketoglutarate and
citrate. Instead of following the canonical reactions of the TCA cycle,
the mitochondrial citrate is exported to the cytosol where it can
participate in non-canonical TCA cycle activity or the malate-aspartate
shuttle to re-enter the mitochondria. Taken together, the dominant carbon
AEFMs predicted by our model suggest HepG2 glutamine is conserved within
TCA intermediates that are regenerated by shuttling mitochondrial citrate
to the cytosol.

\section*{Discussion}

Since the proposal of molecular EFMs, much effort has been made to
enumerate and decompose steady state metabolic fluxes onto these pathways
with little success. By defining EFMs with respect to individual atoms
rather than molecules, we introduced the concept of AEFMs which describe
how individual atoms are remodelled within a metabolic network. As AEFMs
never involve branching pathways, they are easier to interpret than their
molecular EFM counterparts which may involve multiple inputs and outputs
to balance molecular stoichiometries. Using the atom mapping program
RXNMapper, we adapted our previously described CHMC method to efficiently
enumerate carbon and nitrogen AEFMs in five GEMs. Our structural analysis
systematically characterized source metabolite remodelling at the atomic
level. By comparing the number of AEFMs across atoms within the same
source metabolite, we revealed how certain amino acids displayed
remarkable variability in the number of pathways transited by their
constituent atoms. AEFM weight analysis of the HepG2 network revealed that
the majority of glutamine carbons were remodelled via cyclic pathways
regenerating TCA intermediates through cytosolic citrate. Excitingly, our
model also predicted non-canonical TCA activity which has not been
previously described in HepG2 cells.

As a systematic method to explain the flow of metabolites, the problem of
molecular EFM enumeration has historically limited their applicability to
small-scale networks. Many algorithmic advances enumeration have
incrementally improved molecular EFM enumeration, yet never feasibly
beyond networks with $>100$ metabolites and
reactions~\citep{KampAxelvon2006M5fa,TerzerMarco2008Lcoe,vanKlinkenJanBert2016Faet}.
Molecular EFM enumeration is generally regarded as unfeasible for
large-scale networks---so much so that numerous methods have been
developed to enumerate a subset of molecular EFMs constrained by pathway
length~\citep{deFigueiredoLuisF2009Ctse}, reaction
thermodynamics~\citep{GerstlMatthiasP2015Mief}, participating
metabolites~\citep{PeyJon2014Dcoe}, participating
reactions~\citep{KaletaChristoph2009Ctwb}, gene regulatory
rules~\citep{JungreuthmayerChristian2013rSue}, metabolic
subsystems~\citep{SchusterS2002Etps,SchwartzJean-Marc2007Omfa},
extracellular flux
measurements~\citep{SoonsZitaITA2011Iomm,JungersRaphaelM.2011Fcom}, or
even random sampling~\citep{MachadoDaniel2012Rsoe}. While our HepG2
results did show the majority of source metabolite fluxes were explained
by a small number of AEFMs, contextualizing the significance of these
dominant pathway fluxes is only possible when all AEFMs are known.
Furthermore, reaction thermodynamics, gene regulatory rules, and other
constraints may change across different biological conditions. Our results
suggest that enumerating the shortest pathways would fail to recover the
majority of TCA pathways, as the top five ranked AEFMs across all
glutamine carbons averaged 6 reactions in length. By computing AEFMs,
instead, we could feasibly enumerate atomic pathways in genome-scale
networks containing hundreds of metabolites and reactions within minutes
or hours on consumer computing hardware. We are hopeful that algorithmic
tricks borrowed from molecular EFM enumeration tools, such as linear
pathway compression, could further decrease memory usage and run times.

Aside from carbon and nitrogen AEFMs, our work is broadly applicable to
other types of chemical elements. One could feasibly enumerate oxygen or
hydrogen EFMs in small-scale networks, although the significance of these
pathways is unclear for biochemical networks given the ubiquitous nature
of these atoms. We also caution against enumerating hydrogen AEFMs, in
particular, given limitations associated with RXNMapper. Unless explicitly
represented in the reaction SMILES string, RXNMapper does not assign
hydrogen atom mappings. While this may seem trivial, RXNMapper cannot map
atoms in reaction SMILES strings longer than 512 characters. Our method
could also be applied to study sulfur metabolism to quantify methionine or
cysteine remodelling, or phosphorus metabolism to quantify nucleotide
synthesis. Outside the realm of biochemistry, our method is 
applicable to any periodic table element and possibly relevant for
studying general chemical reaction networks. While statistical analyses,
such as different AEFM weight analysis, may reveal metabolic changes in
individual pathways altered across conditions, it is likely that changes
in a metabolic pathway may impact weights in multiple AEFMs. For example,
the majority of AEFMs across distinct source metabolites are not
``molecularly'' independent. By this, we imply that some AEFMs may
correspond to the same sequence of metabolites or different
metabolite/atom sequences stemming from shared reactions. Their weights
will likewise be related through these common reactions. This problem may
perhaps be addressed by formalizing a relationship between atomic and
molecular EFMs. While we have taken significant steps towards solving the
AEFM enumeration and flux decomposition problem, we expect rapid
computational interest analyzing AEFM weights in the near future.

\section*{Materials and Methods}

\subsection*{Reagents and Tools table}

See tables/reagents-and-tools-table.xls.

\subsection*{Methods and Protocols}

\subsubsection*{Datasets}

% Cite remaining 4 datasets here (referenced in the spreadsheet data/tools/software table)
\nocite{OrthJeffreyD2010RaUo}
\nocite{OrthJeffreyD2010RaUoDATASET}
\nocite{BeckerScottA2005Grot}
\nocite{BeckerScottA2005GrotDATASET}
\nocite{BordbarAarash2011iApd}
\nocite{BordbarAarash2011iApdDATASET}
\nocite{ThieleInes2005EMRo}
\nocite{ThieleInes2005EMRoDATASET}

Five highly-curated GEMs were chosen in this study to reflect a diversity
of organisms, metabolites and biochemical reactions. Four networks were
obtained from the BiGG UCSD database~\citep{KingZacharyA2016BMAp} in
ascending order of network size. Datasets were downloaded as SBML (.xml)
files, and the corresponding stoichiometry matrices, metabolites and
reactions were extracted using the Julia SBML.jl package. The HepG2 liver
cancer cell line dataset was chosen for its single set of estimated steady
state fluxes estimated from a combination of protein and metabolite
abundance data. These cells were grown in DMEM with 10\% fetal calf serum,
22\SI{}{\milli\molar} glucose, and 1.8\SI{}{\milli\molar}
glutamine~\citep{NilssonAvlant2020Qaoa}. The steady state fluxes were
recomputed from the flux balance analysis (FBA) model provided by the
authors (figure3\_metabolitebalances.m; MATLAB R2018a)
(\citealt{NilssonAvlant2020Qaoa}; Data ref:
\citealt{NilssonAvlant2020QaoaDATASET}). Across all models, we modelled
all reversible reactions as net forward reactions to simplify the number
of molecular and AEFMs as the majority of experimental and computational
methods can only estimate net fluxes. We assumed the reaction
stoichiometries of the BiGG datasets represented the canonical net
directions of all reversible reaction. For the HepG2 dataset, the net
directions were chosen based on the net fluxes estimated from the FBA
model.

\subsubsection*{Extracting metabolites SMILES strings}

SMILES structures were not encoded in the BiGG SBML files nor the HepG2
dataset. Thus, wee manually extracted SMILES strings corresponding to each
metabolite across all GEMs. Isomeric SMILES strings were taken from
PubChem~\citep{KimSunghwan2023P2u}, the Human Metabolome
Database~\citep{WishartDavidS2022H5tH}, or MetaNetX
database~\citep{MorettiSebastien2021Munf}. Metabolites shared across GEMs
were assigned the same SMILES structure in addition to identical
metabolites stratified across distinct subcellular compartments. Generic
structures such as R-groups were not modelled in this study and were
replaced with a single bonded hydrogen under the assumption this
modification would not alter atom mappings, especially since RXNMapper
does not map hydrogen atoms unless they are explicitly represented in the
SMILES strings. SMILES strings were also not assigned to metabolites with
ambiguous structures nor pseudometabolites with no known chemical
structure. All pseudometabolites and reactions involving pseudometabolites
were removed from the stoichiometry matrix (Table 2).

\subsubsection*{Constructing reaction SMILES strings}

RXNMapper performs atom mapping on a string that encodes the SMILES
structures of all substrates and products for a given reaction. These
reaction SMILES strings are constructed by concatenating SMILES strings of
individual substrates and products. Multiple substrates or products are
separated by the symbol `.', while the reaction arrow delimiting
substrates and products is denoted by the symbol `$>>$'. The substrate and
product SMILES strings are concatenated in order of metabolite appearance
within the stoichiometry matrix, and multiple stoichiometric copies are
consecutively repeated in the reaction SMILES strings. Non-integer
stoichiometric coefficients are therefore not allowed and these
pseudoreactions were removed from the stoichiometry matrix (Table 2). To
maintain steady state flux in the HepG2 network, flux along
pseudoreactions were replaced with unimolecular source and sink fluxes
consuming pseudoreaction substrates and producing pseudoreaction products.
We manually checked several reactions to ensure the assigned atom mappings
were correct. Crucially, we found RXNMapper to fail on nearly all classes
of transaminase reactions that were not ATP-dependent. This problem has
previously been reported in an atom mapping benchmark study prior to the
publication of RXNMapper~\citep{PreciatGonzalezGermanA.2017Ceoa}. These
transaminase reactions were identified in all metabolic networks and
manually atom mapped based on the KEGG reaction mappings (Figure EV6).

\subsubsection*{Pre-processing genome-scale metabolic models}

Several pre-processing steps were required to compute AEFMs from the GEMs.
The stoichiometric coefficients of external reactions were set to one and
their fluxes rescaled to carry equivalent units of source and sink
metabolites. All reactions containing pseudometabolites with no known
SMILES strings were removed from the stoichiometry matrix and replaced
with unimolecular source and sink reactions carrying equivalent units of
flux. The same procedure was also applied to reactions containing
non-integer stoichiometries. These pseudometabolites were next removed
from the stoichiometry matrix rows, and unimolecular reactions involving
the same external or internal metabolites were aggregated. Finally, we had
to remove a small number of reactions whose reaction SMILES strings
exceeded the 512 token limit of RXNMapper. These reactions were likewise
converted into unimolecular source and sink reactions to maintain steady
state flux across the network. We note that all pre-processing steps
described above are wrapped into convenience functions within our Julia
package MarkovWeightedEFMs.jl.

\subsubsection*{Benchmarking enumeration of atomic-versus-standard EFM}

We benchmarked our Julia package MarkovWeightedEFMs.jl
(\textcolor{red}{version X.0.0}) running Julia version 1.10.2 to enumerate
atomic EFMs versus the molecular EFM enumeration program
FluxModeCalculator (\cite{vanKlinkenJanBert2016Faet}) running MATLAB
version R2018a with default parameters including network compression. We
chose FluxModeCalculator because it is the state-of-the-art program for
enumerating molecular EFMs. Benchmarks were performed on a 64-bit Linux
operating system with 64GB of memory and an AMD 5950X CPU with 16 cores/32
threads. All benchmarks reported the wall time to run the atomic or
molecular EFM enumeration functions and did not include startup times to
open Julia/MATLAB, precompile Julia functions, or load packages and input
data. A benchmark time of `did not finish' (DNF) was assigned if either
program could not complete EFM enumeration within 7 days. In practice,
only FluxModeCalculator failed to enumerate the molecular EFMs in three of
the five datasets.

\subsubsection*{Constructing atomic transition graphs}

For a user-specified atom within a given source metabolite, we identify
all possible metabolite/atom states that are traversable by that initial
state and constrained by the reaction atom mapping predictions. For
reactions involving single stoichiometric copies of each substrate and
product, the atom of interest will always move from a single substrate to
product molecule. If this atom is present in a substrate with multiple
stoichiometric copies, we assume there is an equal probability of that
atom occurring in either copy. Depending on the substrate stoichiometric
copy, that atom may move to (i) different product metabolite/atom states,
(ii) the same product but at different atomic positions, or (iii) the same
metabolite/atom position. In all cases, we set the atomic flux of the atom
transitioning from either stoichiometric copy equal to the reaction fluxes
to maintain atomic mass conservation (Expanded Views 5). The resulting
atomic transition graph is completed by connecting any sink fluxes to the
root of the tree initialized on the source metabolite/atom position. 


%As a strongly connected graph, we constructed ACHMCs of each source
%metabolite carbon and nitrogen to enumerate their corresponding AEFMs and
%computed the weights for the HepG2 flux network by steady state analysis
%of these ACHMCs.

\subsubsection*{ACHMC model}

Following \citeauthor{ChitpinJustinG2023AMct}, the atomic transition
graphs were analyzed by an atomic version of our CHMC (ACHMC) method to
enumerate the corresponding AEFMs and compute their weights. In this
formulation, each ACHMC state corresponds to a sequence of metabolite-atom
positions transited by a given source metabolite atom. Each simple cycled
identified by the ACHMC corresponds to an AEFM rooted on a user-specified
source metabolite atom. For the HepG2 flux network, we computed the
probabilities of each AEFM by steady state analysis of the ACHMC. These
AEFM probabilities were scaled to weights such that the totality of
source-to-sink AEFMs was equal to the unimolecular input flux producing
that source metabolite.




\iffalse

\subsubsection*{Model preliminaries}

The structure of a knowledge-driven metabolic flux network is encoded as
an $m$ by $n$ stoichiometry matrix $S$ with $m$ distinct metabolites and
$r$ irreversible reactions. The fluxes along each biochemical reaction $v$ are
positive, real-valued weights that denote the number of times each
reaction is occurring per unit time. Under metabolic steady state, the net
fluxes consuming and producing each metabolite are balanced and $S \cdot
v = 0$. These reactions may be unimolecular--single substrate to single
product transformations-- or involve multiple substrates, products, or
stoichiometric copies of participating metabolites. We do not require the
network to be strong connected; however, our method will treat independent
components separately. In contrast with our previous unimolecular CHMC
model,

we impose additional requirements on the input metabolites and reactions.
The 3D structure of each metabolite must be known and encoded as a SMILES
structure. We recommend, where applicable, using isomeric SMILES
structures to avoid ambiguous or incorrect atom mappings. We further
require integer-valued stoichiometric coefficients because we consider
atoms to be indivisible entities from which we enumerate their AEFMs. We
lastly require the input metabolic flux network to be open such that
external fluxes enter and leave via source and sink metabolites. We note
this is generally not a problem in large-scale metabolic networks; our
requirement is not a technical limitation of our method but enforced to
promote the enumeration of all source-to-sink AEFMs remodelling nutrient
sources.

%Overall, these requirements generally restrict the metabolic flux network
%from containing any pseudometabolites or pseudometabolic reactions (e.g.
%biomass reactions).

A knowledge-driven metabolic network can be decomposed into a finite set
of $K$ molecular EFMs. In a unimolecular reaction network, the metabolites
and steady state fluxes can be represented as a graph and the molecular
EFMs correspond to weighted simple cycles carrying steady state flux
through the metabolite vertices. In a network containing multispecies
reactions, the correspondence between molecular EFMs and simple cycles is
lost as molecular EFMs may include higher order reactions involving flows
across multiple substrates and products. For either type of
knowledge-driven metabolic network, their observed steady state fluxes can
always be represented as a positive, linear combination of their $K$
molecular EFMs. This flux decomposition problem can be represented by the
system of linear equations $Aw = f$, where $A$ is an $r$ by $K$ matrix of
zero or natural numbers denoting the number of times each reaction occurs
molecular EFM $k=1,2,\ldots,K$, where $w_k$ is the molecular EFM weight.
When the number of molecular EFMs exceeds the number of linearly
independent fluxes, $w$ is underdetermined and there are an infinite
number of ways to assign molecular EFM weights to explain the observed,
steady state fluxes.

\subsubsection*{Markovian decomposition of fluxes onto EFM weights} % summary and directives

In our previous work~\citep{ChitpinJustinG2023AMct}, we proposed Markovian
EFM weights as a means to resolve the flux decomposition problem in
networks consisting exclusively of unimolecular reactions. We formulated
a probabilistic model of single particle flux, in the form of a Markov
chain, where we assumed the probability of this hypothetical particle
transitioning from one metabolite state to another was proportional to the
observed, steady state flux. We further proposed a way of efficiently
computing those molecular EFM probabilities and weights by introducing the
\emph{cycle-history Markov chain} which both enumerated the molecular EFMs
and uniquely identified their probabilities which were rescaled to weights
based on the absolute flux magnitudes. To generalize this CHMC method to
networks containing multispecies reactions, we introduce the concept of
AEFMs which can be enumerated and their weights computed from an atomic
variant of our CHMC method (referred to as an atomic cycle-history Markov
chain or ACHMC). The following sections provide intuitive justification of
our method, while we leave our proof of (atomic/single particle) CHMC
algorithm correctness in our original
paper~\citep{ChitpinJustinG2023AMct}.

% should I put the entire proof here? It's a bit too long at 6 pages
% single spaced with 1 inch margins.

\paragraph*{Single atom, discrete-time Markov chain}

\hspace*{0.1cm}\\\noindent

We developed a probabilistic model of atomic flux to compute atomic EFM
weights in steady state metabolic flux networks. At this atomic level,
individual atoms within metabolite undergo biochemical reactions and move
from the substrate to product side of each reaction equation. As the
molecules within metabolic flux network are stoichiometrically balanced
under metabolic steady state, so too are the corresponding atomic fluxes.

For a user-specified atom within a given source metabolite, we identify
all possible metabolite/atom states that are traversable by that initial
state and constrained by the reaction atom mapping predictions. For
reactions involving single stoichiometric copies of each substrate and
product, the atom of interest will always move from a single substrate to
product molecule. If this atom is present in a substrate with multiple
stoichiometric copies, we assume there is an equal probability of that
atom occurring in either copy. Depending on the substrate stoichiometric
copy, that atom may move to (i) different product metabolite/atom states,
(ii) the same product but at different atomic positions, or (iii) the same
metabolite/atom position. In all cases, we set the atomic flux of the atom
transitioning from either stoichiometric copy equal to the reaction fluxes
to maintain atomic mass conservation (Expanded Views 5).

For a given source metabolite/atom state, the subsequent atomic
transitions are encoded as an atomic transition graph $G = (V, E, f)$,
where $V$ is the set of nodes corresponding to
metabolite/atom states, $E$ is the set of directed edges corresponding to
reactions that move the atom from one metabolite to another, and $f \equiv
v$ are the atomic mass fluxes. We assume there exists a special node $v_{ext} \in V$
which represents the external environment. Source flux entering the graph
occurs along the edge $e = (v_{ext}, r)$, where $r \in V - v_{ext}$ is the
source metabolite/atom state. Sink fluxes leaving the graph occur along
edges $e = (i, v_{ext})$, where $i \in V - v_{ext}$ are the sink
metabolite/atom states. The introduction of $v_{ext}$ closes off the graph
and ensures it is strongly connected, where there exists a path in $G$
from any node $i$ to any other node $j$.

We make a Markovian assumption that the movement of an atom from
metabolite/atom position to another in $G$ is proportional to the observed
steady state (atomic) flux of that reaction. We denote these transition
probabilities $T_{ij} = f_{ij} / \sum_{j^{\prime}}f_{ij^{\prime}}$ for $i
\neq j \in |V|$. From here, we follow the same intuition as our
unimolecular CHMC algorithm~\citep{ChitpinJustinG2023AMct} to model atomic
flows within the network. Starting from an initial source metabolite/atom
state $s_0$, the atom of interest will transit a random sequences of
states $s_0, s_1, s_2, \ldots$. As $G$ is strongly connected, that
sequence will, with probability of one, include all states infinitely many
times. If the atom visits state $i$, it will eventually return to state
$i$. The sequence of states in between $s_i, \ldots, s_{i^{\prime}}$
constitutes a cycle, albeit not necessarily a simple cycle. The sequence
many include multiple nested or interleaved simple cycles. However, this
cycle can always be decomposed into a unique set of simple cycles by
repeatedly removing the earliest simple cycle found within the sequence,
which we refer to as the ``first closure reduction rule''. The AEFMs are
the set of all simple cycles in $G$. Thus, we compute the AEFM
weights by setting them proportional to the steady state probabilities of
their corresponding simple cycles.

\paragraph*{Atomic cycle-history Markov chain}

\hspace*{0.1cm}\\\noindent

We compute the steady state probabilities of each AEFM following our
previously described CHMC method for networks of unimolecular
reactions~\citep{ChitpinJustinG2023AMct}. Briefly, the ACHMC is a special
type of Markov chain, where we expand the notion of state to include
sequences of metabolites transited by the atom. The ACHMC therefore traces
the history of a given atom on its way to completing a simple cycle
corresponding to an AEFM.

We first initialize the ACHMC on source metabolite atom $s_0$ within
$G$. Although we could technically choose any arbitrary metabolite/atom
state for $s_0$, it is much easier to interpret AEFMs rooted on
a source metabolite flowing into the network. Every state $S$ in the ACHMC
corresponds to a simple path $s_0, s_1, s_2, \ldots$. If a reaction
transforms $s_m$ to $s_{m+1}$, the corresponding state $S \equiv s_0, s_1,
\ldots, s_m$ transitions to the longer path $S^{\prime} \equiv s_0, s_1,
\ldots, s_m, s_{m+1}$. If a reaction transforms $s_{m+1} = s_i$, the
corresponding state $S \equiv s_0, s_1,\ldots, s_m$ transitions to the
shorter path $S^{\prime} \equiv s_0, s_1, \ldots, s_i$ For either case,
the ACHMC transition probability is $\mathcal{T}_{S,S^{\prime}}
= T_{s_m,s_{m+1}}$.

\paragraph*{Computing steady state AEFM probabilities}

\hspace*{0.1cm}\\\noindent

Once constructed, the ACHMC describes all possible simple cycles that
a given source metabolite atom can transit through. These simple cycle
probabilities are computed by analyzing the steady state dynamics of the
ACHMC. As the atomic transition graph $G$ is strongly connected, the
corresponding CHMC will also be strongly connected. There thus exists
a stationary distribution $\pi$ which corresponds to the steady state
probability of the atom occupying each CHMC state.

To compute AEFM probabilities, we first identify all transitions $E_K$ in
the ACHMC that close simple cycles corresponding to EFM $k$. For
source-to-sink AEFMs, there exists only a single simple cycle and
therefore only one corresponding simple cycle. For internally looped
AEFMs, there may exist multiple simple cycles starting at different
metabolite/atom states, counting towards the same AEFM. We define the
steady state probability of transiting EFM $k$  as the sum of the steady
state probabilities of those transitions or:

$P_k = \sum_{(S,S^{\prime}) \in E_k} \pi_S \cdot
\mathcal{T}_{S,S^{\prime}}$.

\noindent As the AEFM probabilities are proportional to the atomic fluxes,
these probabilities may be scaled to weights based on the absolute
magnitude of the fluxes. Because we enforce that each source metabolite is
produced by an input reaction with stoichiometry of one, the sum of all
corresponding source-to-sink AEFMs must be proportional to the input
metabolite flux. If $l \in K$ denotes all source-to-sink AEFMs
flowing through a given source metabolite $m_0$ and $f_0$ is the input
flux producing that metabolite, we define the proportionality constant
$\alpha = f_0 / \sum_l P_l$. This scaling not only ensures the total
input/output fluxes are correctly reconstructed, but also that every
individual flux in the network is correctly explained by the totality of
AEFM weights.

\paragraph*{Details associated with computing the (A)CHMC stationary distribution}

\hspace*{0.1cm}\\\noindent

Following basic Markov chain theory, the stationary distribution of the
CHMC, $\pi$ is proportional to the eigenvector of the CHMC transition
probability matrix. The eigenvector can be solved by direct or iterative
numerical methods. The choice of algorithm becomes increasingly important
as the size of the CHMC transition matrix grows exponentially in
highly-connected, large-scale networks. We note that iterative eigenvector
solving algorithms from some packages (e.g. Arnoldi method from
Arnoldi.jl, GMRES from IterativeSolvers.jl) resulted in negative
eigenvector coefficients or NaN values. These are possibly numerical
errors and occurred more frequently in larger CHMC matrices. Repeatedly
re-running these methods sometimes resulted in a correct eigenvector,
albeit with significantly longer computation times. We observed that the
linear solver from the Julia package LinearSolve.jl efficiently solved for
the correct eigenvector without error using the default parameters. Hence,
we use and recommend this package to compute the stationary distribution
of large ACHMCs.

\fi

\section*{Data availability}

The code and data to reproduce this study are available at\\
\url{https://www.github.com/jchitpin/reproduce-efm-paper-2024}.

\section*{Acknowledgements}

This research was supported by the Natural Sciences and Engineering Research Council of Canada (NSERC), Discovery grant XXXXX-XXXX-XXXX (RGPIN/06604-2019) and a Digital Research Alliance of Canada (DRAC) Resources for Research Groups grant to TJP (RRG-TPERKINS). J.G.C. was supported by an NSERC CREATE Matrix Metabolomics Scholarship, NSERC Alexander Graham Bell Canada Graduate Scholarship, and Ontario Graduate Scholarship.

\section*{Author contributions}

Justin G. Chitpin: Conceptualization; data curation; methodology; software; formal analysis; investigation; visualization; writing - original draft; writing - review and editing.\\

\noindent Theodore J. Perkins: Conceptualization; methodology; supervision; funding acquisition; investigation; project administration; writing - review and editing.

\section*{Disclosure and competing interests statement}

The authors declare no competing interests.

\bibliography{biblio-aefms}

\section*{Figure legends}

\noindent Fig 1: Atomic cycle-history Markov chain (ACHMC) pipeline. (a) An
example molecular EFM in glycolysis. (b) An example AEFM tracing the flow
of a given glucose carbon in glycolysis. (c) The major steps within our
pipeline to enumerate AEFMs and compute their weights.

\noindent Fig 2: AEFM enumeration is computationally tractable
compared to molecular EFMs in five GEMs. (a) Number of atomic and
molecular EFMs in the five GEMs. (b) The run time of our
MarkovWeightedEFMs.jl method versus FluxModeCalculator running in serial.
(c) Serial run time scaling of MarkovWeightedEFMs.jl as a function of the
number of atomic EFMs. (d) Speedup of selected ACHMCs as a functional of
additional threads.

\noindent Fig 3: Structural analysis of AEFMs in five genome-scale
metabolic models. (a-b) Categories of carbon and nitrogen EFMs. (c-d)
Distribution carbon and nitrogen AEFM lengths. (e-f) Ratio of the greatest
and smallest number of AEFMs between carbons and nitrogens within the same
source metabolite. Representative metabolites from each dataset are
labelled; labels with asterisks indicate that metabolite is also present
in other networks.

\noindent Fig 4: Carbon AEFM weight analysis of the HepG2 metabolic
flux network. (a) A small number of highly active AEFMs explain the
majority of glutamine flow. (b) Subnetwork of the top five AEFMs across
each glutamine carbon collapsed onto metabolites. (c-e) Sankey diagrams of
the top five glutamine carbon AEFMs. (f) Schematic mapping AEFMs onto
traditionally defined metabolic pathways/subsystems. All metabolite names are
labelled following the BiGG database nomenclature.

\section*{Tables and their legends}

\noindent Table 1: Summary statistics for the five metabolic networks in
this study.

\noindent Table 2: Number of metabolites and reactions removed from each
network during pre-processing.

\section*{Expanded View Figure legends}

\noindent Expanded View 1: Network visualization of the five GEMs after
pre-processing. Metabolites (non-grey nodes) are connected to reaction
entities (grey nodes) with node sizes proportional to their degree. Graphs
were constructed in Gephi version 0.10 using the Yifan-Hu proportional
network layout with default parameters. (a) E. coli core. (b) iAB RBC 283.
(c) iIT341. (d) iSB619. (e) HepG2.

\noindent Expanded View 2: FluxModeCalculator scaling across multiple
threads.

\noindent Expanded View 3: Most ACHMCs traverse less than 10\% of the
total metabolite-carbon/nitrogen state space. Each point corresponds to an
ACHMC rooted on a given source metabolite carbon or nitrogen.

\noindent Expanded View 4: Cumulative explained carbon mass flow of input
glucose and amino acid carbons from the HepG2 model. The cumulative
explained mass flow is truncated at 99\% to avoid displaying AEFMs that
carry little atomic mass flow.

\noindent Expanded View 5: Rules for assigning atomic fluxes from
atom-mapped reactions with biological examples. Schematic node colours
represent ground truth atom mappings across substrate(s) and product(s).
Node labels correspond to distinct atomic indices of substrate and product
stoichiometric copies.
(a) For reactions with single stoichiometric copies of substrates and
products, each atom within a given substrate moves to a single atomic
position within a product with atomic flux equal to the reaction flux.
(b-d) For reactions with $s$ stoichiometric copies of a given substrate,
there are $s$ copies of each metabolite-atom state on the substrate and
product side of the reaction. The probability of any atom within
a substrate stoichiometric copy moving to a product metabolite-atom state
remains equal to the reaction flux. (b) The stoichiometric copies of the
substrate fuse together into a single product. (c) The stoichiometric
copies of the substrate form distinct products. (d) The $s$ stoichiometric
copies of the substrate form $s$ copies of products with identical atomic
backbones. (e) Biological example of rule a where the atoms in each
substrate map to a single position on the corresponding product. (f)
Biological example of rule b where two stoichiometric copies of
glutathione fuse together to form a single stoichiometric copy of
glutathione disulfide. (g) Biological example of rule c where two
stoichiometric copies of acetyl-CoA undergo a reaction to produce
a stoichiometric copy of two distinct products.

\noindent Expanded View 6: Examples of incorrectly atom-mapped transaminase
reactions present in one or more of the five networks. (a) Phenylalanine
transaminase. (b) Isoleucine transaminase. (c) Tyrosine transaminase. (d)
Aspartate transaminase.

\end{document}

